\section{Evaluation}
\label{sec:evaluation}

\subsection{Experimental Setup}

We evaluate three distinct questions using archived evidence: whether the
matcher respects its documented policy semantics, whether the existing
verification covers authorization, and what consequences follow from an
explicit authorization divergence in the symbolic model. Implementation
observations and formal results are reported separately. This evaluation
introduces no additional experiments.

The matcher record uses Python~3.14.3 for direct function calls.
The formal records use ProVerif~2.05 on Windows. Three evidence packages
provide the inputs and results: Stage~1 (S), the divergence model (D), and
independent gate controls (G).\footnote{\raggedright Packages under
\texttt{proofs/evidence/}: S = \texttt{authz-stage1-\allowbreak
20260916T123834944479Z/}; D = \texttt{authz-turepass-\allowbreak
20260919T090826295518Z/}; G = \texttt{authz-chat-gates-\allowbreak
20260918T023621017067Z/}.}
Their manifests record execution metadata, source revisions, commands,
input hashes, and output provenance; raw standard output, standard error,
and verification summaries accompany the runs. The reported model inputs
match their recorded hashes. These archives originate from separate
revisions, rather than a single integrated service experiment.

Table~\ref{tab:evaluation-runs} reports the six communication-model runs
used here. Each completed with exit status zero. Counts refer to final
query results, not to satisfied security properties: reachability is
reported through negated-event queries. Durations are single recorded
verifier runs, rounded to one decimal place. They do not measure service
overhead or establish a performance comparison; repeated measurements and
a controlled hardware benchmark are outside this evaluation.

\begin{table}[htbp]
\centering
\small
\begin{tabular}{llrr}
\hline
Package & Model / scenario & Results & Seconds \\
\hline
S & Original communication & 11 & 334.6 \\
S & Authorization diagnostic & 12 & 353.7 \\
D & Divergence (\texttt{turepass}) & 2 & 598.4 \\
G & Provider allow / Receiver allow & 5 & 813.6 \\
G & Provider deny / Receiver allow & 5 & 507.1 \\
G & Provider allow / Receiver deny & 5 & 552.5 \\
\hline
\end{tabular}
\caption{Archived ProVerif runs used in the evaluation.}
\label{tab:evaluation-runs}
\end{table}

% S: agent_communication.pv; agent_communication_authz.pv.
% D: agent_communication_authz_chat_turepass.pv.
% G: agent_communication_authz_chat_p_allow_r_allow.pv;
% agent_communication_authz_chat_p_deny_r_allow.pv;
% agent_communication_authz_chat_p_allow_r_deny.pv.
% Source: the three package manifests and verification-summary.txt files.

\subsection{Matcher Evaluation}

The intended semantics select the most specific matching rule. The
self-contained \texttt{proofs/test\_matcher\_bug.py} illustrates the
inconsistency with a specific deny rule for
\texttt{mallory@example.com:*} followed by a wildcard allow rule.
This explanatory script defines its own matcher variants and simplified
specificity function; it does not import the implementation matcher.

Package S supplies direct implementation evidence in
\texttt{policy-\allowbreak matcher-\allowbreak sanity.json}.
Its recorded invocation uses
\texttt{alice@example.com:agent}, a specific rule
\texttt{alice@example.com:*} with budget $-1$, and a wildcard rule with
budget $10$. The rulebook is valid, and the recorded specificity scores
are $70$ and $0$. The documented semantics therefore require budget $-1$
in either ordering. Table~\ref{tab:matcher-evaluation} reports the saved
results for these two orderings.

\begin{table}[htbp]
\centering
\small
\begin{tabular}{lrr}
\hline
Rule order & Expected budget & Returned budget \\
\hline
Specific deny, then wildcard allow & $-1$ (deny) & $10$ (allow) \\
Wildcard allow, then specific deny & $-1$ (deny) & $-1$ (deny) \\
\hline
\end{tabular}
\caption{Direct matcher observations for the recorded Alice rulebook.}
\label{tab:matcher-evaluation}
\end{table}

The first ordering establishes a concrete policy decision mismatch; the
second shows that reordering the same rules changes the result. This is
an implementation-level, order-dependent authorization bypass. The record
covers matcher calls, not Provider processing, token issuance, or message
acceptance in a running service. Two observed orderings do not estimate
the prevalence of the defect across rulebooks or deployments.

\subsection{Formal Verification Results}

\paragraph{Authorization coverage gap.}
The original communication model in S produces eleven results: seven
event-reachability results, three authentication correspondences, and one
token-secrecy result. The diagnostic model preserves these results and
adds the authorization correspondence
\[
\mathsf{Accept}(R,I,t)\Longrightarrow\mathsf{PolicyAllow}(R,I).
\]
This added correspondence is false, while the three authentication
correspondences and token secrecy remain true. Here, \texttt{Accept}
denotes receipt and verification of a token by the initiating peer;
it is not application-message acceptance.

The diagnostic process never emits \texttt{PolicyAllow}. Consequently,
the failed correspondence exposes an authorization obligation absent from
the original verification; it does not independently demonstrate a
matcher execution error. The retained authentication and secrecy results
apply to these Stage~1 queries. They are not automatically inherited by
the later communication models.

\paragraph{Authorization consequence model.}
Package D contains two final results for \texttt{turepass}.
The query asserting that \texttt{ChatAccept} is unreachable is false,
with a reconstructed formal reachability witness. Under preset
\texttt{ScenarioSpecDeny} and \texttt{ScenarioImplAllow} facts, the
witness contains intended denial, material release, token issuance, and
application-message acceptance. This is a symbolic reachability
consequence of the divergence abstraction.

The second result establishes
\[
\mathsf{ChatAccept}(R,I,t,m)\Rightarrow\mathsf{TokenIssue}(R,I,t).
\]
Thus, acceptance can coexist with the assumed policy denial while still
having corresponding prior token issuance. Credential provenance and
policy permission remain separate obligations. Issuance is observed in
the witness; this model has no standalone \texttt{TokenIssue} reachability
query and no \texttt{ChatAccept}--\texttt{ChatSend} correspondence query.
The complete D archive supplies these two final results with independent
execution provenance, beyond the historical \texttt{turetrace} excerpt.

\paragraph{Authorization gate experiments.}
Package G evaluates three fixed pairs of Provider and receiver permissions.
Each control has three reachability queries and two correspondences.
Table~\ref{tab:gate-evaluation} translates the reachability results into
reachable (R) or unreachable (U), avoiding ambiguity about negated-query
truth values.

\begin{table}[htbp]
\centering
\small
\begin{tabular}{llccc}
\hline
Provider & Receiver & Release & Issue & Accept \\
\hline
Allow & Allow & R & R & R \\
Deny & Allow & U & U & U \\
Allow & Deny & R & U & U \\
\hline
\end{tabular}
\caption{Gate-control reachability: \texttt{ProviderRelease},
\texttt{TokenIssue}, and \texttt{ChatAccept}.}
\label{tab:gate-evaluation}
\end{table}

Provider denial prevents release and downstream actions. Receiver denial
prevents issuance and acceptance, while Provider release remains reachable.
The allow/allow control establishes an executable permitted path.
Both acceptance-to-issuance and acceptance-to-send correspondences are
true in all three controls. In the denial cases they hold vacuously,
because acceptance is unreachable; only the allow/allow case supplies
a non-vacuous control. These experiments locate blocking stages within
the model, rather than evaluate a deployed repair. The controls use a
fixed private token, unlike the fresh tokens in D, and are not a verified
refinement of the divergence model.

\subsection{Security Findings Summary}

The evidence supports three complementary findings. First, direct matcher
calls demonstrate a decision inconsistent with documented specificity
semantics. Second, the coverage diagnostic distinguishes authorization
from the authentication and secrecy properties actually queried in
Stage~1. Third, given explicit divergence assumptions, the consequence
model permits message acceptance, while the gate controls identify
distinct release and issuance boundaries.

Together, these findings support the SAGA case study's authorization
enforcement gap. The implementation observation establishes the local
defect; the formal results characterize modeled consequences and control
positions. Neither evidence layer substitutes for the other. No result
establishes cryptographic primitive failure, certificate forgery, key
compromise, or arbitrary attacker control of the modeled peer processes.

\subsection{Limitations}

The formal model abstracts the matcher result. It neither executes Python
nor evaluates the concrete rulebook, so the implementation record and the
formal reachability witness are not a single executable reproduction.
A strict semantic mapping from concrete rules, specificity, and returned
budgets to scenario facts and formal events remains necessary.

Private transport and the \texttt{ProviderReleased} handoff are model
abstractions. The evaluated queries do not establish token expiry, quota
consumption, revocation, policy-update behavior, session isolation, or
injective agreement. No complete production-service reproduction or
deployment-scale measurement is reported. Future work may extend the
semantic mapping and evaluate service behavior and larger deployments;
the present results establish neither their outcomes nor generality
across agentic systems.
